% Condensed review article prepared for Current Science (ISSN 0011-3891).
% Source of truth for this submission.  Formatted by build_cs_docx.py into
% Current_Science_Manuscript.docx following the official Instructions to Authors
% and Guidelines for Manuscript Preparation.
%
% Deliberate editorial decisions, recorded for the author:
%   * body condensed from ~9000 to <6000 words (journal limit for Review Articles)
%   * figures reduced from 10 to 6; detailed tables moved to supplementary material
%   * first-person narration removed (the journal requires a formal, objective tone)
%   * citations renumbered serially by first appearance and set as superscripts
%   * British spelling throughout

\title{Calibrated set-valued prediction: a review of conformal methods}
\shorttitle{Conformal prediction review}

\begin{abstract}
A confidence score is not a measured frequency. Conformal prediction replaces a
point prediction with a set calibrated on held-out data, containing the true label
with probability at least $1-\alpha$ under exchangeability alone. This review
organises two decades of work by the guarantee pursued: marginal validity,
sharpness, risk control and conditional validity. A reproducible study on a
ten-class benchmark gives 0.8996 coverage at a 0.90 target, shows confidence
thresholds falling to the model's 0.8207 accuracy, and reports online adaptation
recovering 0.9111 coverage under shift at a large cost in set size. Validity is
cheap marginally and expensive conditionally.
\end{abstract}

\begin{keywords}
calibration; conformal prediction; distribution-free inference; risk control;
uncertainty quantification
\end{keywords}

\section{Introduction}

Automated decision systems usually return a single answer and a number that looks
like a probability. In image classification that number is the softmax score, which
is the output of a model trained to minimise cross-entropy, not a measured
frequency. A model can be systematically overconfident on data unlike its training
set, and a threshold chosen on one sample gives no assurance on the next. This
matters wherever an error rate is specified in advance: a screening pipeline that
must miss no more than one case in twenty, or a diagnostic aid whose output is
acted on by a clinician.

Conformal prediction, introduced by Vovk, Gammerman and Shafer \cite{Vovk2005} and
made accessible by Shafer and Vovk \cite{Shafer2008}, addresses this by changing
what is calibrated. Instead of asking the model to produce honest probabilities, a
separate calibration sample is used to calibrate the model's decisions. For each
calibration example a nonconformity score measures how poorly the observed label
agrees with the model's output; an empirical quantile of these scores defines which
labels are admitted to the prediction set for a new input. The resulting set
satisfies

\begin{equation}
\Pr\{Y_{\text{new}} \in C(X_{\text{new}})\} \geq 1-\alpha
\end{equation}

for any joint distribution of the data, provided only that calibration and test
examples are exchangeable. No assumption is made about the model, the distribution
of the covariates, or the relation between them. The model may be badly specified
and frequently wrong; the guarantee still holds. What is not guaranteed is that the
set will be small. The tension between validity and efficiency is the main subject
of the field's subsequent history.

Three developments turned a theoretical construction into a usable method. A
single split of the data suffices, so the guarantee costs one model fit and one
pass over the calibration sample \cite{Lei2018}. Scores can be constructed so that
set size responds to the difficulty of the individual input rather than to the
model's average behaviour \cite{Romano2020,RAPS}. And the guarantee
extends from coverage to any bounded loss, so the same calibration machinery can
control false-negative rates and other quantities that applications actually care
about \cite{Bates2021,LTT}. Running alongside these results is a strong
negative one: distribution-free conditional coverage, which would bound the error
rate separately for every subgroup or every value of the covariates, cannot be
achieved in any non-trivial sense \cite{Barber2021}. Much current work consists of
finding principled ways to live with that impossibility.

This review is organised by guarantee rather than by algorithm. Each generation
reuses the same three stages of fitting, calibrating and thresholding; what changes
is the object being calibrated and the assumption that the calibration relies on.
The reader is assumed to be a scientist outside machine learning, so technical
terms are defined where they first appear. The literature search was carried out on
4 November 2026 using OpenAlex and the arXiv API. Every entry in the reference list
carries a DOI or an arXiv identifier resolved during that search; preprints are
identified as such in the text. Chinese-language journals are not covered because
CNKI and Wanfang do not expose open interfaces, and this gap is stated rather than
filled with unverified citations.

Three surveys are close enough to require a distinction. The tutorial of
Angelopoulos and Bates \cite{Angelopoulos2021intro} establishes the intuition of
split conformal prediction for work up to 2022 and is assumed here. Campos and
colleagues \cite{Campos2024} survey natural language processing, organising by
task. Zhou and colleagues \cite{Zhou2025} take a data perspective, organising by
modality and by the assumption each modality violates. The guarantee-based axis
used here is complementary to the data axis; the two disagree about which work is
central, and those disagreements are noted below.

Figure 1 gives the chronology. The foundational period is unexpectedly long. From
2005 to about 2018 the field possessed a valid and general construction but saw
little practical adoption; what made adoption possible was the recognition that
one split suffices. Work density rises sharply after 2021, concentrated in
conditional guarantees and language-model applications.

\section{What the guarantee assumes}

Exchangeability means that the joint distribution of the data is unchanged by any
permutation of the indices. It is weaker than the familiar assumption of
independent and identically distributed observations, and every independent and
identically distributed sample satisfies it. Any ordering that carries information
breaks it: time series, active learning, distribution drift that depends on the
order in which examples arrive, and any pipeline in which the calibration sample
was selected using information about the test points.

A calibration sample drawn at random from the same collection as the test set is
exchangeable with it. A calibration sample drawn from last year's data and deployed
against this year's is not. In the latter case the guarantee does not degrade
gracefully; it fails silently. Coverage on the deployment distribution can fall far
below $1-\alpha$ while every diagnostic computed on the calibration distribution
looks correct. Section 7 measures this failure mode directly.

Figure 2 shows the construction. A model is fitted on a training sample. A disjoint
calibration sample of size $n$ is passed through it to obtain nonconformity scores
$s_i=s(X_i,Y_i)$, where a larger score means that the pair agrees less well with
the model. The scores are sorted and the quantile $q$ is taken at index
$\lceil (n+1)(1-\alpha)\rceil$. For a new input, the prediction set contains every
label whose score does not exceed $q$.

The finite-sample correction in that index, $(n+1)$ in place of $n$, is what makes
the guarantee exact rather than asymptotic, and it is the detail most often
implemented incorrectly. It buys validity for every sample size at the price of one
effective calibration observation. The implementation used in Section 7 applies
this correction and, when a threshold would produce an empty set, forces the
argmax label into the set; this can only raise coverage and therefore preserves
the lower bound.

Two properties of the construction are easily confused. The guarantee is
marginal: it averages over the randomness of the calibration and test samples and
says nothing about a particular test point. A procedure covering 90 per cent of
cases on average may cover 99 per cent of one subgroup and 60 per cent of another
without violating it. The guarantee is also unconditional on the model. A better
model yields smaller sets, a worse model yields larger sets, and neither yields an
invalid one. This is why conformal prediction is described as a wrapper around a
model, and why research on efficiency concentrates on the score rather than on the
validity theorem.

\section{Scores and the efficiency question}

Once the quantile fixes validity, all design freedom lies in the score. The
simplest classification score is $s(x,y)=1-\hat p_y(x)$, the complement of the
model's estimated probability for the true class. This produces the set of all
labels whose predicted probability exceeds a global threshold. Its weakness
appears as soon as there are many classes: a class that the model systematically
underestimates is excluded more often than the nominal error rate allows, and
coverage is recovered only by compensating across classes.

Adaptive scores instead accumulate probability mass in descending order until the
true label is reached, so the score depends on where the true class sits in the
ranking. The adaptive prediction set (APS) score adds the cumulative probability
of all classes ranked above the true class to a random fraction of the true class's
own probability \cite{Romano2020}; the randomisation restores exact validity when
there are ties. A regularised variant penalises the number of classes in the set
and suppresses large sets at a small cost in coverage \cite{RAPS}. In
regression the corresponding choice is the absolute residual, possibly scaled by
an estimate of local variance, and the prediction set becomes an interval
\cite{Lei2018}.

There is no distribution-free answer to the efficiency question, and this is a
theorem rather than a gap in the literature. The expected size of any conformal
prediction set is bounded below by a quantity depending on the Bayes risk of the
problem; no score can produce sets smaller than the ambiguity intrinsic to the
data. Scores can approach that bound, and the empirical literature is largely a
sequence of increasingly clever attempts to do so.

\section{Marginal versus conditional coverage}

Marginal coverage averages over the covariate distribution. Conditional coverage
would require

\begin{equation}
\Pr\{Y \in C(X) \mid X = x\} \geq 1-\alpha \quad \text{for almost every } x,
\end{equation}

which is a statement about each point rather than about the population as a whole.
Barber and colleagues \cite{Barber2021} proved that no distribution-free procedure
can satisfy this except by returning the trivial set of all labels. The intuition
is that a finite calibration sample contains only finite information about any
particular $x$, and a procedure required to be valid for every distribution can
always be defeated by one that concentrates its mass elsewhere.

The practical consequence is that work on conditional coverage does not attempt
the impossible. It either weakens the target to validity averaged within known
groups, or adds assumptions under which the target becomes reachable. Figure 3
illustrates both the schematic point and its measured form. In the experiments of
Section 7, the naive score attains marginal coverage 0.8996 against a nominal
0.90, while its class-conditional coverage ranges from 0.776 to 0.995 across the
ten digits. Both numbers are correct. Reporting only the first gives a true
statement that misleads about the second, and the gap is a property of the
construction rather than an error of implementation.

\section{Four generations of methods}

The four generations are defined by the guarantee each pursues; the ordering is
one of dependence rather than strict chronology, because some third-generation
ideas predate the maturation of second-generation methods. Figure 4 summarises the
genealogy.

The first generation established that distribution-free marginal validity is
attainable and cheap. The original transductive construction tested each candidate
label by examining how well the augmented sample conformed, at the cost of
refitting the model for every candidate \cite{Vovk2005}. Lei and colleagues showed
that a single split into training and calibration portions preserves the coverage
guarantee at a fraction of the cost \cite{Lei2018}. Sadinle and colleagues
formalised the classification case as set-valued classifiers with bounded error
levels and introduced the criterion of least ambiguity that later efficiency work
adopted as its objective \cite{Sadinle2019}. What this generation could not do was
make set size respond to the difficulty of an individual input. With a fixed score
and a global threshold, size is determined by the model's average behaviour.

The second generation took validity as given and competed on set size. Adaptive
scores made size respond to the shape of the predicted distribution rather than to
its maximum. Tibshirani and colleagues extended the guarantee to covariate shift
by reweighting calibration scores according to the likelihood ratio between the
test and calibration covariate distributions \cite{Tibshirani2019}; validity is
preserved under the specified shift, at the cost of requiring that likelihood
ratio to be known or estimable. The remaining limitation is the one identified in
Section 4: sharpness averaged over the population says nothing about sharpness for
the subgroup that matters.

The third generation moved from coverage to arbitrary bounded risk. In image
segmentation the quantity of interest may be the fraction of mislabelled pixels; in
screening it may be the false-negative rate. Bates and colleagues proposed
risk-controlling prediction sets that bound the expectation of a monotone loss
\cite{Bates2021}. The Learn then Test framework subsequently reframed calibration
as multiple hypothesis testing, controlling risk with finite-sample validity
\cite{LTT}. The same development that made the machinery general also
clarified its limits, through the impossibility result described above.

The fourth generation accepts that impossibility and works around it. It either
targets a relaxed form of conditional validity under explicit assumptions, relaxes
exchangeability itself, or moves to settings where the loss is not a simple
indicator of coverage. Table 1 compares the four generations along the dimensions
that decide whether a method can be used in a given application.

\section{Recent developments, 2024--2026}

Two constructive routes now exist around the impossibility result. One assumes
smoothness: if the conditional coverage function varies slowly with $x$, its local
estimate can be calibrated, giving a guarantee that is conditional up to a
smoothness constant. Hore and Barber \cite{Hore2024} showed that randomisation is
what makes local weighting identifiable, and obtained robust guarantees under weak
conditions. The other route targets groups. Gibbs, Cherian and Candès
\cite{Gibbs2025} require coverage to hold on average within each of a known set of
groups, which is weaker than pointwise conditional coverage but needs no
smoothness assumption. This is the natural formalisation of a fairness requirement
for diagnostic systems. Each route makes its extra assumption explicit, which is
an improvement over the pre-2021 situation.

Two recent surveys reorganise the field along axes that the guarantee-based view
does not capture. Campos and colleagues \cite{Campos2024} catalogue which
constructions survive the violations of exchangeability typical of natural
language processing, where the predicted object is a sequence and the loss is
often not coverage. Zhou and colleagues \cite{Zhou2025} organise by data modality
and by the assumption each modality violates. Work central on one axis can be
peripheral on the other, and reading both gives a more accurate map than either
alone.

Much of the recent activity concerns large language models, for structural reasons.
Their outputs are sequences, so the natural prediction set is a set of sequences,
which is infinite and cannot be enumerated; the loss that matters is usually
factual error rather than label coverage; and deployment in which the model is
updated, the prompt distribution drifts, or the retrieval corpus changes violates
exchangeability. One line of work defines nonconformity scores over sequences and
calibrates a stopping rule \cite{ConformalLM}. A second moves to risk control,
bounding the expected fraction of unsupported claims \cite{CRC}. A third treats
prediction sets as diagnostics: Alur and Das \cite{Alur2026} observe that set size
is itself a signal of distribution shift, making it a cheap, model-independent
monitor for data corruption.

A number of 2026 preprints extend these directions to hierarchical group risk
control, introspective scores for vision-language models, paraphrase-invariant
scoring and the auditing of validity under model change. They are not cited
formally here because they have not completed peer review; their claims should be
read with that qualification. Other recent work extends conformal prediction to
generative output spaces under covariate shift \cite{Li2026gen}, to ordinal labels
where the loss respects distance \cite{OrdinalCRC}, and to cross-validated
calibration that avoids the efficiency loss of a single split \cite{XVCrc}.

\section{An original empirical study}

A survey that reports only what others measured inherits their choices about what
was worth measuring. This section reports a small, fully reproducible study
designed to make three things visible: that the coverage guarantee holds for a
severely undertrained model, that confidence thresholding cannot replace it, and
that validity is not maintained for free under distribution shift. The study is
deliberately small so that every number can be checked; it makes no claim about
scale.

The data are the scikit-learn digits benchmark: 1797 examples, 64 features and ten
classes, with pixel values scaled to the unit interval. Ten random seeds each
produce a stratified split into a training pool of 898, a calibration sample of
449 and a test sample of 449. The model is a standardised multinomial logistic
regression with $C=1.0$ and 3000 iterations. Only 5 per cent of the training pool,
about 44 examples, is used to fit it, giving a test accuracy of
$0.8207\pm0.0236$. This label-scarce regime is a deliberate choice. With the full
training pool the accuracy exceeds 0.97, the conformal quantile becomes so small
that almost every prediction set collapses to a singleton, and the
coverage--efficiency trade-off disappears from view. It is also the regime in which
conformal prediction is most useful in practice, since a well-trained model on
abundant data needs the guarantee less. All reported quantities are means and
standard deviations over the ten splits. Table 2 gives the settings; the full
sweep table is in the supplementary material.

Coverage and efficiency across six values of $\alpha$ appear in Figure 5. At
$\alpha=0.10$ the naive score reaches $0.8996\pm0.0189$ against a target of 0.90,
a deviation of 0.0004, with mean set size $1.331\pm0.164$. The guarantee holds
where it should. At $\alpha=0.30$ the target is 0.70, the model's accuracy is
0.8207 and the measured coverage is 0.8207 with mean set size 1.000: the
singleton already clears the target, so the procedure is valid for a reason
unconnected with conformal calibration. The APS score is systematically
conservative at intermediate $\alpha$, exceeding the target by 0.044 at
$\alpha=0.10$; this is a known property of the randomised cumulative score at
moderate sample sizes and is the price of a flatter class-conditional profile.

The second experiment asks whether the guarantee can be obtained more cheaply by
thresholding the model's own confidence. Four procedures are compared at a nominal
0.90 target: the raw softmax threshold, the same threshold after temperature
scaling fitted by minimising calibration negative log-likelihood, split conformal
with the naive score, and split conformal with APS. The results are in Table 3. The
fitted temperature averaged $0.749\pm0.088$, so the rescaling was substantial, yet
the first two rows are identical in coverage and set size. Temperature scaling is
a monotone transformation of the logits; it preserves the ranking of classes and
therefore the argmax, so a threshold on rescaled probabilities selects the same
singleton set as a threshold on the original ones. Calibrating a score, in the
sense of making probabilities match frequencies, is a different task from
calibrating a decision, in the sense of making the set's error rate match a
target. A paired Wilcoxon signed-rank test over the ten seeds gives a statistic of
4.0 ($p=0.0137$) for the comparison of absolute coverage deviation between the
raw threshold and APS, a result compatible with a real but modest effect.

The third experiment violates exchangeability in a controlled way. A stream of 450
test examples is divided into nine consecutive blocks, and Gaussian noise with
standard deviation increasing from 0 to 0.60 is added to the pixels of each block
before clipping to the unit interval. A static procedure calibrates once on clean
data and keeps the resulting quantile. An adaptive procedure starts from the same
quantile but updates $\alpha$ after each observation according to

\begin{equation}
\alpha_{t+1} = \alpha_t + \gamma\,(\alpha_t - 1\{Y_t \notin C_t\}),
\end{equation}

with $\gamma=0.02$ and clipping to $[0.001,0.999]$, following Gibbs and Candès
\cite{Gibbs2021ACI}. Figure 6 shows the outcome. The static procedure attains 0.4644
coverage over the whole stream, falling from 0.90 in the clean block to 0.24 at a
noise level of 0.45. The adaptive procedure attains 0.9111 overall, with
block-wise coverage between 0.76 and 0.98, despite receiving no information about
the noise level. Its mean set size rises from 1.256 to 7.618 classes. In a
ten-class problem this approaches the trivial set: the guarantee is real, but the
information carried by the set is largely spent. The final block's coverage of
0.76 reflects the step size, which bounds how quickly $\alpha_t$ can move; the
limitation belongs to the update rule rather than to the framework, and choosing
$\gamma$ remains an unresolved trade-off between responsiveness and stability.

The study establishes four points. Marginal validity holds as stated for a model
of 0.82 accuracy. Confidence thresholding, including after temperature scaling,
does not achieve it. Class-conditional coverage can deviate widely while marginal
coverage meets its target. Static calibration loses validity under shift, and
online adaptation recovers it at a large efficiency cost. The study does not
establish that conformal prediction works at scale, that these numbers transfer to
full-size image benchmarks, or that the adaptive procedure is preferable in any
specific deployment. The dataset is small, the model is deliberately weak, and the
corruption is synthetic and one-dimensional.

\section{Open problems and outlook}

Theoretical obstructions and engineering gaps should be kept apart. Among the
former, every construction since the impossibility result buys conditional
validity with an assumption, and the assumptions are not interchangeable:
smoothness, known groups and estimable likelihood ratios lead to guarantees of
different strength with different failure modes. A unified account of which
assumption buys which guarantee, at what efficiency cost, does not yet exist.
Barber and colleagues \cite{Beyond} give coverage guarantees under a
quantified departure from exchangeability, but the quantification is a
total-variation-style quantity that cannot usually be estimated from data;
restating the guarantee in terms of an observable diagnostic would be more useful.
For language models, the natural object is a distribution over sequences, and
sequence-level coverage is either trivial or undefined; whether a useful analogue
exists remains open.

Among engineering gaps, calibration data are treated as free although in practice
they are expensive, and the sample-size requirement interacts with the target:
the finite-sample correction consumes one effective observation, so at
$\alpha=0.01$ with a small calibration set the quantile is dominated by a few
order statistics. Composite systems are a second gap: a conformal set passed to a
downstream model gives that model a set rather than a point as input, and
exchangeability of the components does not imply exchangeability of the whole.
Third, the observation that set size can serve as a monitor is immediately
deployable but underdeveloped; a monitoring procedure based on it would need a
null distribution, a detection threshold and a false-alarm guarantee.

Finally, conformal prediction and selective inference share a methodological
commitment: finite-sample guarantees obtained by calibrating against held-out data
rather than by assuming a parametric model. Stability selection and knockoffs
control false discovery by constructing comparison distributions
\cite{Meinshausen2010,Barber2015,Candès2018}, while conformal prediction controls
coverage by constructing a quantile. Recent work on conditional calibration of
knockoffs makes the connection explicit \cite{Luo2025}. A general theory of
calibration-based inference covering both appears within reach.

\section{Conclusions}

Conformal prediction delivers what it promises: a finite-sample, distribution-free
guarantee of marginal coverage, obtained by calibrating a model that need not
itself be calibrated, correct or even good. Work since 2023 has extended the
promise in two directions, from coverage to arbitrary bounded risk and from
marginal to conditional guarantees under explicit assumptions, while also
establishing where unconditional conditional validity is impossible. The
unresolved question is the practical cost of the guarantee. Validity under shift
is recovered only by inflating sets until they carry little information,
conditional validity requires assumptions that are often unverifiable, and the
calibration data on which the whole construction rests are scarce rather than free.

For practice, three statements suffice. If the requirement is a declared error
rate on exchangeable data, split conformal with the finite-sample quantile
correction delivers it and confidence thresholding does not, however well the
probabilities are calibrated. If the requirement is error rates within named
subgroups, group-conditional calibration can deliver it when the groups are known.
Under drift, static calibration fails silently; online adaptation restores
validity, but the resulting set sizes must be measured before deployment rather
than discovered afterwards.

\begin{acknowledgements}
The computations were carried out with scikit-learn, NumPy, SciPy and Matplotlib.
No external funding was received. The author thanks the anonymous reviewers for
comments that improved the presentation.
\end{acknowledgements}

\begin{supptables}
Supplementary material contains the complete coverage--efficiency sweep over six
values of $\alpha$ for both scores, the full experimental settings table, the
2024--2026 contributions table with publication status for each entry, four
additional figures, and the analysis script together with the JSON file from which
every reported number is read.
\end{supptables}
